% Experiment report. Fixed sections: see experiments/TEMPLATE/report.tex and experiments/README.md.
% Build: npm run experiment:report -- 30-rl-army-level
\documentclass{../polyexp}
\begin{document}
% generated by scripts/experiment-report.ts from experiment.json and models/<run>/ - do not edit
\providecommand{\ExpHeader}{}\renewcommand{\ExpHeader}{%
{\sffamily\footnotesize\color{expgrey}EXPERIMENT 30 \textperiodcentered\ METHOD \hfill 2026-11-28}\par
{\LARGE\bfseries\sffamily A lasting army level in the state, value iteration on r15}\hfill\verdictbox{expred}{NOT KEPT}\par\smallskip
{\itshape With the army factor (train and attack) as a lasting level in the state, does value iteration find when to build and use more or less army, and does that beat r15?}\par\medskip
\begin{tabular}{@{}p{22mm}p{\dimexpr\linewidth-26mm}@{}}
\fact{Change}{Experiment 29 with the action set army5 instead of econ5: the base bot's train and attack factors (move types 12 and 2, together, as spend3's army) at x0.5, x0.75, x1, x1.33 or x2, held between turns and moved one step at most per turn, from x1; value iteration over (turn, level) and (turn, level, income bucket).}
\fact{Knobs}{\knob{state} $\in$ \{turn,level, turn,level,income/3\}; \knob{minCount} $\in$ \{20, 100, 500\}, best 500}
\fact{Base}{\model{greedy-full-probe-perf-r15} $\rightarrow$ result \model{rl-gate vs greedy-full-probe-perf-r15: (turn, level) FAIL +7, with income FAIL +14, static army x0.5 FAIL +10; the army factor barely changes r15's games (77\% unchanged at x0.75 / x1.33)}}
\fact{Method}{300k games with a random walk over the levels; value iteration for both state versions and min-count 20, 100, 500; min-count chosen per version on 2,000 validation seeds; both chosen policies on the 10,000 evaluation seeds against r15, and against the best fixed level of the action check if one beats r15; metric: final raw score, \#171 rule}
\fact{Setups}{Perfection, 16$\times$16 continents, 4 players, Crazy, seat 0 Xin-Xi (human seat), learner base greedy-full-probe-perf-r15; opponents greedy-full-probe-perf-r15, greedy-full-probe-perf, greedy-full-probe-perf (anchor), seats shuffled per seed (rl-gate defaults from 27 Nov)}
\fact{Record}{issue \issue{166}; status: done}
\end{tabular}\par\medskip
\begin{center}\begin{tikzpicture}[node distance=5mm]
    \node[flowbox] (n0) {{\sffamily\bfseries\scriptsize\color{expgrey}BASE}\\[1pt]\model{full-probe-perf-r15}};
    \node[flowbox,right=5mm of n0] (n1) {{\sffamily\bfseries\scriptsize\color{expgrey}CHANGE}\\[1pt]\knob{state}\\{\scriptsize turn,level, turn,level,income/3}\\\knob{minCount}\\{\scriptsize 20, 100, 500}};
    \node[flowbox,right=5mm of n1] (n2) {{\sffamily\bfseries\scriptsize\color{expgrey}MEASURED}\\[1pt]Paired, 10,000 evaluation seeds\\{\scriptsize 1 setup, metric: final raw score, \#171 rule}};
    \node[flowbox,right=5mm of n2] (n3) {{\sffamily\bfseries\scriptsize\color{expgrey}RESULT}\\[1pt]Mean final score, 10,000 evaluation seeds: r15 vs army-level policy (turn, level)\\{\scriptsize 7199.0 $\rightarrow$ 7206.0}\\\textbf{+7.0}};
    \node[flowbox,right=5mm of n3] (n4) {{\sffamily\bfseries\scriptsize\color{expgrey}VERDICT}\\[1pt]\textbf{\color{expred}NOT KEPT}\\{\scriptsize new champion:}\\\model{rl-gate vs greedy-full-probe-perf-r15: (turn, level) FAIL +7, with income FAIL +14, static army x0.5 FAIL +10; the army factor barely changes r15's games (77\% unchanged at x0.75 / x1.33)}};
    \draw[flowarrow] (n0) -- (n1); \draw[flowarrow] (n1) -- (n2); \draw[flowarrow] (n2) -- (n3); \draw[flowarrow] (n3) -- (n4);
\end{tikzpicture}\end{center}
\begin{center}\sffamily\small\begin{tabular}{@{}lrrr@{}}\toprule
Measure & Base & New & Change\\\midrule
Mean final score, 10,000 evaluation seeds: r15 vs army-level policy (turn, level) & 7199.0 & 7206.0 & +7.0\\
Same seeds: r15 vs army-level policy (turn, level, income) & 7199.0 & 7213.0 & +14.0\\
\bottomrule\end{tabular}\end{center}
}
\providecommand{\ExpResults}{}\renewcommand{\ExpResults}{%
{\small\color{expgrey}No runs recorded.}
}

\ExpHeader

\section{Question and hypothesis}
\textbf{The idea (David, 28 Nov 2026, issue \#166).} Experiment 29 for army: the train and attack factors (move types 12 and 2, together, as spend3's army) becomes a lasting
level in the state, which the policy moves one step per turn. In 29 the economy level passed against r15 ($+195$),
but the gain came from lowering economy overall, not from its timing. This experiment asks the same question for
army, with economy at r15's weights.

\textbf{Prediction} (before any run of this experiment):
\begin{enumerate}
  \item \emph{No fixed level helps much.} In the action check, no level played every turn beats r15 by more than 100
        points ($\times2$ every turn is the same as spend3's army action, $+9$ in 25 and 26).
  \item \emph{Little to gain.} The chosen policy's mean difference against r15 on the evaluation seeds is between
        $-50$ and $+100$, and it does \emph{not} pass the \#171 rule.
  \item \emph{Less army late.} Units trained in the last turns rarely pay back in Perfection: without income in the
        state, the path is at $\times0.75$ or below from turn 25 on.
  \item \emph{The model is optimistic, as in 29:} its predicted gain exceeds the measured gain by more than 200
        points.
  \item \emph{Income adds little}: within $\pm 60$ of the version without, not significant.
\end{enumerate}
\textbf{Against the prediction:} a PASS would mean that r15's army weight is off for Perfection, as its economy and
research weights were, or that the timing pays.

\section{Mechanism}
As experiment 29, with the action set \texttt{army5} (\texttt{rl\_value::strategic}): the base bot's train and attack factors (move types 12 and 2, together, as spend3's army)
times 0.5, 0.75, 1, 1.33 or 2. The levels are symmetric around $\times1$ because the direction is not known in
advance (29's economy levels went down to $\times0.25$ because the check had shown that direction). At $\times1$ the
weights are r15's, bit for bit (a test checks). The level starts at $\times1$ and moves at most one step per learner
turn; a state without a policy entry keeps the level.

\textbf{State}: (turn, level before the turn) and (turn, level, income / 3). \textbf{Value iteration}: \texttt{rl-value
fit-level}, certainty-equivalence model, $\gamma = 1$, one backward sweep, min-count $m$ (see 29). \textbf{Data}: a
random walk over the levels from $\times1$ (\texttt{--driver walk}).

\section{Setup}
As experiment 29: learner base \texttt{greedy-full-probe-perf-r15}; opponents r15, probe-perf and the anchor
probe-perf, shuffled per seed; Perfection, 16$\times$16 Continents, 4 players, Crazy, learner Xin-xi in the human
seat. Seeds as in 29: training $1\dots300{,}000$, action check $600{,}000{,}000\dots$ (2{,}000), validation
$700{,}000{,}000\dots$ (2{,}000), evaluation $10^9\dots$ (10{,}000).

Commands (from \texttt{engine/}):
\begin{verbatim}
B=./target/release/rl-value; R=target/rl-value; E=../experiments/30-rl-army-level
$B check --actions army5 --games 2000 --seed-start 600000000 --out $E/check.json
$B collect --actions army5 --driver walk --games 300000 --seed-start 1 --out $R/army5-walk.tsv
$B fit-level --train $R/army5-walk.tsv --income-bucket 0 --min-count 20,100,500 --out $R/level-army5
$B fit-level --train $R/army5-walk.tsv --income-bucket 3 --min-count 20,100,500 --out $R/level-army5
$B eval --actions army5 --policy $R/level-army5/policy-bB-mM.tsv --n 2000 --seed-start 700000000 --out $E/val-bB-mM.json
rl-gate --candidate rl-policy:experiments/30-rl-army-level/policy-bB-mM.tsv --base greedy-full-probe-perf-r15
\end{verbatim}

\section{Results}
\ExpResults

All numbers are copied from \texttt{check.json}, \texttt{level-summary-b*.json}, \texttt{val-*.json},
\texttt{eval-*.json} and the \texttt{rl-gate} JSONs in this folder (28 Nov 2026; per-seed rows dropped, the runs are
deterministic). Data: 300{,}000 games with a random walk over the levels on r15, 8{,}512{,}460 rows, 1{,}482\,s on
10 threads. States: 144 without income, 1{,}750 with income / 3.

\paragraph{Action check} (2{,}000 check seeds, each level every turn; r15 7{,}148): $\times0.5$ 7{,}196 ($+48$,
532 better / 493 worse), $\times0.75$ 7{,}149 ($+1$, 222 / 241), $\times1.33$ 7{,}149 ($+1$, 205 / 234), $\times2$
7{,}157 ($+9$, 418 / 427). Most games do not change at all: at $\times0.75$ and $\times1.33$ about 77\% of the seeds
end with exactly r15's score.

\paragraph{Value iteration and validation} (2{,}000 validation seeds against r15, base mean 7{,}282).

\begin{center}\small
\begin{tabular}{lrrrrrr}
\toprule
version & $m$ & $V(\text{start})$ & keep $\times1$ & predicted gain & validation mean $d$ & better / worse \\
\midrule
(turn, level)          & 20, 100, 500 & 6{,}767 & 6{,}687 & $+80$  & $+4$  & 126 / 131 \\
(turn, level, income)  & 20           & 6{,}938 & 6{,}685 & $+253$ & $-21$ & 192 / 207 \\
                       & 100          & 6{,}934 & 6{,}685 & $+249$ & $-22$ & 190 / 207 \\
                       & 500          & 6{,}906 & 6{,}685 & $+221$ & $-20$ & 200 / 209 \\
\bottomrule
\end{tabular}
\end{center}

Chosen: $m = 500$ for both versions (the three are the same policy without income; the best of three, all within
noise, with income).

\paragraph{Evaluation seeds} (10{,}000, \texttt{rl-gate}; r15 7{,}199, median 8{,}655, p10 1{,}480):

\begin{center}\small
\begin{tabular}{llrrrrl}
\toprule
candidate & base & mean & mean $d$ (99\% bounds) & better / worse & sign $p$ & verdict \\
\midrule
(turn, level)         & r15 & 7{,}206 & $+7$ ($-8$, $+22$)   & 748 / 715     & 0.20  & FAIL \\
(turn, level, income) & r15 & 7{,}213 & $+14$ ($-3$, $+32$)  & 1{,}084 / 1{,}017 & 0.075 & FAIL \\
static $\times0.5$ & r15 & 7{,}209 & $+10$ ($-20$, $+40$) & 2{,}698 / 2{,}594 & 0.078 & FAIL \\
\bottomrule
\end{tabular}
\end{center}

Equal seeds: 8{,}537 and 7{,}899 of 10{,}000 (correlation 0.98 with r15). p10 1{,}490 and 1{,}485. 15 opponents
(information only): $+44$ (616 / 488, $p = 7\cdot10^{-5}$) and $-14$ (348 / 301).

\section{Behaviour}
Mean army factor by turn on the evaluation seeds (\texttt{eval-b*-m500.json}):

\begin{center}\small
\begin{tabular}{lcccccc}
\toprule
version & turns 0--4 & 5--9 & 10--14 & 15--19 & 20--24 & 25--29 \\
\midrule
(turn, level)         & 0.70 & 0.70 & 0.90 & 1.02 & 1.20 & 1.08 \\
(turn, level, income) & 0.78 & 1.20 & 1.13 & 0.85 & 0.90 & 0.93 \\
\bottomrule
\end{tabular}
\end{center}

Without income the path goes down to $\times0.5$ in turns 4--6, back to $\times1$ by turn 9, and to $\times1.33$ on
6 of turns 18--27; overall $\times0.5$ / $\times0.75$ / $\times1$ / $\times1.33$ on 11 / 35 / 36 / 18\% of turns, mean
factor 0.93. With income the policy goes the other way in the opening (up to $\times1.56$ at turn 8) and stays near
$\times0.9$ later; mean factor 0.97. The two versions disagree on the direction in almost every phase, a sign that
the values of the levels are about equal.

\section{Interpretation}
\textbf{Prediction 1 (no fixed level beats r15 by more than 100): right} (best $\times0.5$, $+48$).
\textbf{Prediction 2 ($-50$ to $+100$, no PASS): right} ($+7$ and $+14$, both FAIL).
\textbf{Prediction 3 ($\le\times0.75$ from turn 25): wrong.} The path is at $\times1$ or $\times1.33$ in turns 25--28
and at $\times0.75$ only in turn 29.
\textbf{Prediction 4 (model optimistic by more than 200): half right.} With income it predicts $+221$ and gains $+14$;
without income it predicts only $+80$ and gains $+7$.
\textbf{Prediction 5 (income adds little): right} ($+7$ between the two means).

\textbf{The army weight hardly matters to r15 in this setup.} Scaling the train and attack factors by 0.75 or 1.33
leaves three quarters of the games unchanged, and halving or doubling them moves the mean by less than 50. The
bot's army moves are decided by the terms that do not scale (threat and capture terms, unit costs), or the
alternatives are far apart in score, so a factor of two rarely flips the choice. With so little signal, value
iteration picks among nearly equal levels, and the path is noise: it zigzags, and the version with income chooses
the opposite direction in the opening. Unlike economy (29) and research (31), r15's army weight is not off for
Perfection.

\textbf{The model's optimism follows the number of states.} Without income (144 states) the predicted gain is $+80$,
with income (1{,}750 states) $+221$; the measured gain is about zero in both. More states mean more maxima over
nearly equal, noisy $Q$ values, and so more maximisation bias.

\section{Threats to validity}
\begin{itemize}
  \item \emph{One coupling of train and attack.} Both factors move together, as spend3's army did. Training more
        while attacking less (or the reverse) is not tested and could matter where the sum does not.
  \item \emph{Opponents and setup.} Against 3 Crazy bots in Perfection, fights decide less than in Domination; with
        15 opponents the $\times$-level policy without income gains $+44$ (information only), so army may matter
        more in crowded games.
  \item \emph{Factor range.} $\times0.5$ to $\times2$ may be too narrow to change the bot's army moves often; a
        wider range (or a lower floor) was not tried.
\end{itemize}

\section{Verdict and next step}
\textbf{Not kept}: both policies fail the \#171 rule against r15 ($+7$, $+14$, sign test not significant). A
negative result: r15's army weight is about right for Perfection, and the timing of army does not pay either.

The best fixed level of the action check, $\times0.5$ (\texttt{schedule-static-x0.5.tsv}), fails too ($+10$).

Next step (one idea): split train and attack, or test the army level in the 15-opponent setup, where the info run
hints at a small gain.

\section{Notes for the paper}
Draft claim: the same level method that finds a $+195$ economy and a $+152$ research gain finds nothing for army,
because the army weight barely changes the bot's play (77\% of games unchanged at $\pm25\%$). A useful control:
value iteration does not invent gains where the data has none, while its model still predicts $+80$ to $+221$.
Carried by the action check and the evaluation table, next to 29 and 31.
\todo{jakob}{Is this worth a place in the paper? Where?}

\end{document}
